% !PRINTED 98
\weberpage{98}
%<<CONT
wähnten Stelle gezeigt ist. Wir wollen hier die Rechnung in etwas
veränderter Form anordnen. Wir setzen:

\[ \xi = \eta - \frac{\alpha}{3}, \]

was zur Folge hat, dass $(\varrho, \eta) = (\varrho, \xi)$ ist, wegen
$1 + \varrho + \varrho^2 = 0$. Dann ergiebt sich für $\xi$ die cubische
Gleichung:

\weq{16}{\xi^3 + P\xi - Q = 0,}

worin

\weq{17}{P = \beta - \frac{\alpha^2}{3}, \quad Q = \frac{2\alpha^3}{27} - \frac{\alpha\beta}{3} + \gamma.}

Nun ergiebt sich aus (12):

\[ \begin{aligned}
(\varrho, \xi)(\varrho^{-1}, \xi) = m &= \xi^2 + \xi'^2 + \xi''^2 - \xi\xi' - \xi\xi'' - \xi'\xi'' \\
&= -3P
\end{aligned} \]

oder

\weq{18}{P = -\frac{m}{3},}

und aus (13), nach einigen einfachen Rechnungen, wenn

\[ S = \xi^2\xi' + \xi'^2\xi'' + \xi''^2\xi, \quad S' = \xi\xi'^2 + \xi'\xi''^2 + \xi''\xi^2 \]

gesetzt wird:

\[ \begin{aligned}
m(a + b\varrho) &= 9Q + 3\varrho S + 3\varrho^2 S' \\
m(a + b\varrho^2) &= 9Q + 3\varrho^2 S + 3\varrho S' \\
0 &= 9Q + 3S + 3S',
\end{aligned} \]

woraus durch Addition:

\weq{19}{27Q = m(2a - b),}

und durch Subtraction der beiden ersten:

\weq{20}{mb = 3(S - S') = 3(\xi - \xi')(\xi'' - \xi)(\xi'' - \xi').}

Wenn $m$ durch 9 theilbar ist, so ist $\alpha = 0$, und $P$ und $Q$ sind
\gesperrt{ganze} Zahlen und $P$ durch 3 theilbar. Ist aber $m$ nicht durch 9 theilbar,
so sind erst $3P$ und $27Q$ ganze, nicht durch 3 theilbare Zahlen.

Ist $D$ die Discriminante der Gleichung (15), also eine ganze Zahl, so ist
auch

\weq{21}{\begin{aligned}
\sqrt{D} &= (\eta - \eta')(\eta'' - \eta)(\eta'' - \eta') = (\xi - \xi')(\xi'' - \xi)(\xi'' - \xi') \\
&= \frac{mb}{3}
\end{aligned}}

eine ganze Zahl, und also ist, wenn $m$ nicht durch 9 theilbar ist, $b$
durch 3 theilbar. Nun ist nach (14)
%CONT>>
